% Root and B full literal audits passed. No numerical table premise is used.
\subsection{A ninth excess with a proper strong pair}

\begin{proposition}\label{C:strong-pair-excess-nine-through21}
Let $C$ be a connected finite category with $k$ objects and $n$ arrows.
Suppose no object is a retract of another, $C$ has a proper strongly
connected component on two objects, and
\[
 n\leq\min\{21,2k+8\}.
\]
Then $BC$ is not homotopy equivalent to $\mathbb CP^2$.
\end{proposition}
\begin{proof}
The cases with $n\leq16$ follow from Theorem~\ref{C:CP2-lower17}.
Write $S$ for the strong pair, $s=|S|$, and $\ell=k-2\geq1$.
The no-retract hypothesis puts both identities outside its minimal
ideal, so $s\geq6$. Contracting $S$ in the undirected support gives
at least $\ell$ edges, and at least one interface expands into
arrows incident with both objects of $S$. Counting the outside
identities therefore gives
\[
 n\geq s+2\ell+1,\qquad n\leq2\ell+12,\qquad s\leq11.
\]
For $s\leq10$, the categorical Rees comparison gives
$H^2(BS;\mathbb Q)=0$: nonzero second homology requires total
row and column counts at least three, hence at least nine ideal
arrows and the two additional identities.

\paragraph{An extreme pair.}
By opposition suppose $S$ is a source. Let the complement $D$ have
$q$ connected components, and let $T$ be the total number of arrows
from $S$ to $D$. Every component receives at least two such arrows.

For $s=6,7$, the minimal group is trivial with one row and one
column at each object. The actual source-block collapse of
Lemma~\ref{C:one-row-source-block-collapse} saves $s-1$ internal
arrows and at least one interface arrow. Its result has at most
$n-s\leq15$ arrows, which is impossible.

Suppose $8\leq s\leq10$. If $H^2(BD;\mathbb Q)=0$, the
two-block cup lemma finishes. Otherwise the excess-two lemma gives
\[
 |D|\geq2\ell-q+3,\qquad
 n\geq s+2\ell+q+3\geq s+2\ell+4.
\]
The upper bound forces $s=8$, $q=1$, $|D|=2\ell+2$ and $T=2$.
There is exactly one external arrow from each object of $S$, and
all other external hom-sets at $S$ are empty. The minimal group of
$S$ is trivial, and its total row/column counts are $(2,2)$,
$(3,2)$ or $(2,3)$. It has one minimal row per object or one
minimal column per object. All external entrances and exits are
singleton or empty, and hence fixed by its minimal idempotents.
Lemma~\ref{C:one-row-fixed-entrance-collapse}, or its opposite,
replaces $S$ by a point through an actual homotopy equivalence,
without increasing any other hom-set. At least seven arrows
disappear, leaving at most fourteen.

For $s=11$, the elementary connectivity bound
\[
 n\geq11+(2\ell-q)+2q=2\ell+11+q
\]
forces $q=1$, $|D|=2\ell-1$ and $T=2$.
Thus $D$ has trivial corners and singleton hom-sets on an
undirected tree; its classifying space is contractible.
Both interface arrows end at a single object $w$ of $D$,
because an interface meets both objects of $S$. There is no
cross arrow out of $w$ in $D$, since composing it with the two
interface arrows would contradict $T=2$. The full subcategory
on $S\cup\{w\}$ has terminal object $w$.
The nerve of $C$ is exactly the union of that contractible
subcategory nerve and $ND$, meeting along the point $w$.
This is a union of CW subcomplexes and is contractible.
This excludes every extreme case.

\paragraph{A nonextreme pair.}
Choose outside objects $a,b$ with $a\to S\to b$.
Composition supplies $a\to b$. The contracted undirected
support therefore has at least $\ell+1$ edges, including two
interfaces. Let $P_{\rm ext}$ be the parallel excess summed
over all individual hom-sets between $S$ and its complement.
The count is
\[
 n\geq s+2\ell+3+P_{\rm ext},\qquad
 P_{\rm ext}\leq9-s,\qquad s\leq9.
\]

For $s=6,7$, the unique minimal-ideal endomorphism $e_i$ at
each object is a zero of its corner. There are four possible
one-sided fixed-arrow conditions: all entrances or all exits
at either object are fixed by its $e_i$. Internal extra arrows
can violate at most $2(s-6)$ of these conditions, and external
parallel excess can violate at most $P_{\rm ext}$ more.
This is the same literal four-condition count proved in
Proposition~\ref{C:strong-pair-excess-eight-through21}.
Consequently a condition holds whenever
\[
 2(s-6)+P_{\rm ext}<4.
\]
The one-sided deletion lemma applies. Its unit hypothesis holds
because a unit fixes every arrow fixed by a zero idempotent.
The deleted object removes at least its two corner arrows,
the two internal cross arrows, and an external entrance and
exit. The result has at most $n-6\leq15$ arrows.
This treats $s=6$ and $s=7$ with $P_{\rm ext}\leq1$.

For $s=8,9$ and $P_{\rm ext}=0$, the minimal group is trivial,
and the minimal band has one row per object or one column per
object, as its size is at most seven. All external hom-sets
are singleton or empty. The fixed-entrance block collapse or
its dual applies as above, saving at least seven arrows.

The only remaining numerical possibilities are
\[
 (s,P_{\rm ext})=(7,2),\qquad(8,1).
\]
In either case equality holds throughout the lower-bound count:
$n=2\ell+12$, all outside corners are trivial, every outside
cross hom-set is a singleton, the contracted support has exactly
$\ell+1$ edges, and exactly two of those edges meet $S$.
Its unique cycle is the triangle $a,S,b$. All other contracted
edges form trees attached to that triangle. They cannot attach
directly to $S$, since there are only two interfaces.

Every nonempty such tree has a leaf outside $\{a,S,b\}$.
It is an object with trivial corner and a single singleton
hom-set incident with another outside object of trivial corner.
The singleton cofinal deletion removes it by an actual homotopy
equivalence. Repeating deletes all $\ell-2$ extra outside objects,
two arrows each time. The four-object category that remains has
\[
 n-2(\ell-2)=16
\]
arrows, and is excluded by Theorem~\ref{C:CP2-lower17}.
If $\ell=2$ there was no tree to delete and the same final bound
holds already. This completes the nonextreme cases.
\end{proof}

\begin{corollary}\label{C:eighteen-five-acyclic-scope}
A category of type $\mathbb CP^2$ with eighteen
arrows and at least four objects has either four objects, or
five objects and acyclic support. In the five-object case its
corner, parallel and support-cycle excesses satisfy
$D+P+R=9$.
\end{corollary}
\begin{proof}
The accepted object bound gives at most five objects. A proper
retract deletion would leave at most fifteen arrows, so no
retract exists. Proper strong components on at least three
objects are excluded through twenty-one arrows by
Proposition~\ref{C:large-SCC-minimal-candidate21}, and a wholly
strongly connected category needs at least twenty-three by
Corollary~\ref{C:strong-category-CP2-minimum23}.
For five objects the preceding proposition excludes every
proper strong pair, because $18=2\cdot5+8$. Acyclicity follows.
The excess identity is $18-(2\cdot5-1)=9$.
\end{proof}
